\section{Safety Benchmarks}

安全评估不是简单的``会不会拒答''。安全意味着很多东西：有害内容、越狱、双重用途、上下文相关的风险，以及不同地区法律与社会规范的差异。

\subsection{HarmBench 与 AIR-Bench}

HarmBench 关注违反法律或规范的有害行为；AIR-Bench 则从监管框架和公司政策出发，把风险做了更细的分类。它们都说明：安全不是单一指标，而是一组情境化的约束。

\begin{table}[H]
\centering
\small
\begin{tabular}{p{2.5cm}p{3.0cm}p{3.5cm}p{4.0cm}}
\toprule
\textbf{Benchmark} & \textbf{关注点} & \textbf{输入类型} & \textbf{评估方式} \\
\midrule
HarmBench & 有害内容生成 & 对抗性 prompt（手写 + 自动生成） & 分类器判断是否成功引出有害内容 \\
AIR-Bench & 监管与政策合规 & 基于真实法规分类的测试用例 & 按政策类别逐项评估 \\
TruthfulQA & 事实幻觉 & 容易引出常见错误的问题 & 人工标注的真实性判断 \\
BBQ & 社会偏见 & 歧义/消歧义的选择题 & 统计偏差方向和幅度 \\
\bottomrule
\end{tabular}
\caption{主要安全 benchmark 的对比：侧重点不同，适用场景也不同}
\end{table}

\subsection{Jailbreaking}

GCG 之类的方法可以自动优化 prompt，诱导模型绕过安全策略。这个现象说明：安全评估不能只依赖模型表面的拒答行为，还要考虑其是否真的抵抗攻击。

越狱攻击的演进大致分为几个阶段：
\begin{itemize}
    \item \textbf{手工 prompt 注入}：早期通过角色扮演（``DAN''）、前缀注入等方式绕过安全策略。
    \item \textbf{自动化搜索（GCG）}：Zou et al. 提出了 Greedy Coordinate Gradient 方法，自动搜索能绕过安全策略的对抗性后缀。
    \item \textbf{多模态攻击}：通过图片、音频等非文本通道注入指令，绕过纯文本层面的安全过滤。
    \item \textbf{多轮渐进式攻击}：不在一个 prompt 中直接要求有害内容，而是通过多轮对话逐步引导模型偏离安全策略。
\end{itemize}

\begin{warningbox}{安全评估是动态攻防，不是静态打分}
与知识类 benchmark 不同，安全评估面对的是主动对手。攻击手段在不断演进，因此``通过了当前的安全评估''不等于``安全''。模型发布后，必须持续监测新出现的攻击方式，并定期更新评估标准。这也是为什么部分安全研究机构主张在模型上线前做``红队压力测试''，而不是仅靠固定的 benchmark 套件。
\end{warningbox}

\subsection{Safety 评估中的 taxonomy 问题}

安全评估最困难的部分之一是``什么算有害''本身没有统一定义。不同文化、法律体系和使用场景对``有害''的界定差异巨大：

\begin{itemize}
    \item \textbf{法律差异}：某些内容在一个国家合法但在另一个国家违法（如仇恨言论在美国受第一修正案保护，但在德国是违法的）。
    \item \textbf{上下文依赖}：``如何制作炸弹''在化学教育语境和犯罪语境中完全不同。
    \item \textbf{程度问题}：从轻微不礼貌到严重的物理伤害指导，``有害''是一个连续谱而非二分类。
    \item \textbf{时间变化}：社会规范在演变，三年前被认为可以接受的回答，今天可能被视为不当。
\end{itemize}

\begin{knowledgebox}{为什么安全 benchmark 需要定期重做而不只是扩充}
与知识 benchmark 不同，安全 benchmark 面临的不仅是``题目泄露''问题，还有``定义漂移''问题。随着社会规范和法规的变化，什么算``安全''的标准本身在移动。因此，安全评估套件不能只扩充题目，还需要定期审核和更新分类体系（taxonomy）本身。
\end{knowledgebox}

\subsection{Pre-deployment testing}

课程提到美英安全研究机构在模型发布前参与测试。这里的关键不是某个具体流程，而是一个原则：在模型进入公众之前，先做更系统的安全审查。

\subsection{什么是安全}

安全既包含 \textbf{capability}，也包含 \textbf{propensity}。

\begin{itemize}
    \item 一个模型可能有能力做坏事，但平时会拒绝。
    \item API 模型里，propensity 特别重要。
    \item 开源模型里，capability 更重要，因为安全层很容易被微调移除。
\end{itemize}

\begin{warningbox}{双重用途}
CyBench 之类的任务既可以被看作安全评估，也可以被看作能力评估。一个系统越会做攻防，越需要明确它是在帮助防御，还是在增强攻击能力。
\end{warningbox}

